%% ============================================================
%%  Appendix C — Extensions: Traffic Realism, Escalation
%%  Tracking, and Dynamic Response
%% ============================================================
\chapter{Extensions: Traffic Realism, Escalation Tracking, and Dynamic Response}
\label{app:extensions}

This appendix documents a subsequent round of practical extensions to the
SEDM implementation described in Chapters~3--5: more realistic synthetic
traffic generation, a severity-weighted alternative to the sliding-window
escalation signal, and two non-learning mechanisms for adjusting honeypot
behaviour dynamically in response to observed traffic. It closes with a
revisiting of the learning-based alternative discussed in
\S\ref{sec:dqn_discussion}, in light of this project's own history with
the DQN baseline.

\section{Motivation and Scope}

Three practical gaps remained in the implementation evaluated in
Chapter~4. First, synthetic traffic was generated by sampling every
network feature independently at every simulation step from a fixed
per-attack-type distribution; no session had a persistent ``character,''
and all benign traffic was drawn from one distribution regardless of
context. Second, the escalation-rate signal feeding the SEDM's R3
override and the composite risk score was a hard sliding window over
binary attack/no-attack outcomes, discarding information about attack
severity. Third, the SEDM's decision matrix and its thresholds were
entirely static; the system had no mechanism for remembering a source IP
across sessions or for adjusting its own behaviour in response to
sustained observed patterns.

A fourth question --- whether a learning-based approach such as the DQN
baseline from Chapter~3 would be a more practical way to address the
third gap --- is addressed directly in \S\ref{sec:app_c_dqn}, using
evidence already established by this thesis's own DQN results
(\S5.2) and the evaluation-harness defect history documented in
Appendix~B and the project's bug log.

\section{Synthetic Traffic Realism}

\subsection{Session-coherent feature generation}

The \texttt{Attacker} agent now draws a persistent \textbf{session
profile} once per episode: an intensity scalar
$\iota \sim \text{Lognormal}(0, 0.35)$ applied multiplicatively to eight
volume-shaped network features (\texttt{sbytes}, \texttt{dbytes},
\texttt{sload}, \texttt{dload}, \texttt{spkts}, \texttt{dpkts},
\texttt{ct\_srv\_src}, \texttt{ct\_dst\_ltm}), and a benign-traffic
persona selection (\S\ref{subsec:personas}). TTL, window size, duration,
and packet-loss features are left unscaled, as session intensity has no
clear physical interpretation for those fields. This is a deliberately
minimal, additive change: \texttt{\_simulate\_features()} gained
keyword-only \texttt{intensity}/\texttt{persona} override parameters
defaulting to the session's own profile, so a whole episode reads as one
coherent attacker machine rather than independently-sampled noise at
every step, while any caller needing genuinely independent samples
(\S\ref{subsec:retraining}) can still request them explicitly.

\subsection{Benign traffic personas}
\label{subsec:personas}

Three benign-traffic personas replace the single fixed \texttt{NORMAL}
distribution: \texttt{casual\_user} (the original baseline, weight 0.70),
\texttt{crawler} (short, high-volume, single-service requests, weight
0.20), and \texttt{monitoring\_probe} (tiny payloads, many distinct
destinations, weight 0.10). Persona selection never affects
\texttt{attack\_type} or kill-chain state --- only which of the fifteen
raw feature values is sampled --- but materially affects the diversity of
benign traffic the classifier and policy are evaluated against, which
matters directly for the false-positive characterisation discussed in
Finding~3 of \S5.2.

\subsection{Training-data variance and classifier retraining}
\label{subsec:retraining}

Two call sites --- \texttt{AttackClassifier.generate\_training\_data()}
and the feature-distribution visualisation in \texttt{main.py} ---
construct one \texttt{Attacker} per class and draw many samples from it.
Without an explicit per-sample override, every ``independent'' training
sample of a class would silently share one session's intensity/persona,
understating the true feature variance the classifier needs to learn
from. Both call sites were updated to draw a fresh intensity/persona per
sample from a local random number generator. The shipped classifier was
retrained on the resulting distribution; holdout accuracy on a
freshly-seeded batch was 99.4\%, consistent with the near-perfect
separability already reported in Finding~3 of \S5.2.

\subsection{Randomised event payloads}

Independently, the OpenCanary event emulator's payload templates --- a
small, literal, enumerable set of example credentials and paths per
scenario --- were parameterised with wordlists and jitter ranges resolved
at generation time, so repeated invocations of the same scenario produce
varied-but-realistic payloads rather than one of a handful of exact
strings. This affects only the \texttt{logdata} field of emitted events;
the kill-chain classification pipeline reads only the numeric
\texttt{logtype} code and is unaffected.

\section{Escalation Tracking: From Hard Window to Severity-Weighted EMA}

The original escalation-rate signal,
\begin{equation}
  \rho_{\text{window}} = \frac{1}{|W|}\sum_{i \in W} \mathbb{1}[\text{attack}_i],
\end{equation}
over a fixed window $W$ of the most recent 20 steps, is binary per step
and has a hard cutoff: a step that falls outside the window stops
influencing the rate immediately rather than fading out, and a
reconnaissance probe counts identically to a worm outbreak. A
severity-weighted exponential moving average was added as an opt-in
alternative:
\begin{equation}
  \rho_{\text{ema}}^{(t)} = \alpha \cdot s(\text{attack}_t) + (1-\alpha)\,\rho_{\text{ema}}^{(t-1)},
  \qquad \alpha = 0.15,
\end{equation}
where $s(\cdot)$ is the existing per-attack-type severity weight from
Table~\ref{tab:full_hyperparams} (Appendix~B). This is $O(1)$ per step
rather than requiring deque maintenance, decays smoothly, and reflects
attack severity rather than mere frequency. Because it is
severity-weighted, its steady-state ceiling under sustained
maximal-severity activity is bounded by the highest severity weight
(0.90, for WORMS) rather than 1.0 --- a calibration difference from the
window signal that means \texttt{RATE\_THRESHOLD} (\S3.4, R3) would need
independent re-validation before EMA mode is relied upon operationally.
For this reason \texttt{"window"} remains the default in both the
training environment and the live pipeline; both signals are computed
and exposed regardless of which one is selected, so the two can be
compared directly.

The identical formula is implemented in the live OpenCanary pipeline's
\texttt{SessionTracker}, so the escalation signal is computed
consistently regardless of which of the two parallel data paths
(synthetic training environment or live event stream) produced it.

\section{Dynamic Response: Cross-Session Reputation and Bounded Threshold Adaptation}

\subsection{Cross-session reputation and the R4 override}

Session state in the live pipeline is scoped to one source IP for the
duration of its TTL (300\,s by default) and is then discarded; a source
IP that attacks, falls silent past the TTL, and returns is treated as
entirely new. A \texttt{ReputationTracker} component was added to address
this: a per-source-IP offense score, stored independently of individual
sessions, that decays with a configurable half-life (default six hours)
rather than expiring at a fixed time:
\begin{equation}
  \text{score}(t) = \text{score}(t_0) \cdot 0.5^{\,(t - t_0)/h}.
\end{equation}
Every event increments the reputation score for its source IP by
\texttt{offense\_increment} $\times$ severity (0.25 by default), clamped
to 1.0. This feeds a new fourth override rule, \textbf{R4}, evaluated
before R1: if a source IP's decayed reputation meets or exceeds
\texttt{REPUTATION\_THRESHOLD} (0.60), the action is escalated one
severity level regardless of how benign the immediate event appears.
Positioning R4 ahead of R1 is a deliberate security-posture decision --- a
source IP that has already accumulated sufficient reputation does not
regain unconditional trust merely by sending one innocuous-looking packet
--- with the corresponding and openly acknowledged trade-off that a
shared or dynamically-reassigned IP address that was briefly malicious
remains penalised for the reputation half-life even after any actually
malicious activity from that address has ceased.

\subsection{Bounded adaptation of the R3 threshold}

A second mechanism, \texttt{AdaptiveThresholds}, allows
\texttt{RATE\_THRESHOLD} --- the trigger for the R3 override --- to drift
within a bounded range in response to its own observed trigger frequency.
Its scope was chosen deliberately narrowly.
\texttt{ESC\_LOW\_THRESHOLD} and \texttt{ESC\_HIGH\_THRESHOLD}, which
bucket the analytically-computed escalation risk from the Markov
transition model (\S3.4), were excluded from any adaptive mechanism:
because that risk is computed purely from the static, hand-specified
transition matrices, no live observational signal exists in this system
whose drift would indicate the two band cuts require adjustment, and
constructing an adaptive mechanism against a fixed, non-observational
input would not constitute a genuine improvement. \texttt{RATE\_THRESHOLD},
by contrast, gates on the empirically observed escalation rate --- but
the only signal actually available for it is R3's own firing frequency,
not whether any individual trigger was operationally correct, since no
ground-truth feedback pipeline exists in this system
(\S\ref{sec:app_c_dqn}). \texttt{AdaptiveThresholds} is therefore scoped
and documented honestly as an alert-fatigue safety valve --- keeping R3's
trigger rate near an operator-specified target so that a downstream alert
queue neither floods nor goes silent --- rather than as a correctness
improvement, and implemented as a bounded deadband controller: once an
observation window (200 decisions) has accumulated, the threshold is
nudged by a fixed step if the observed R3 rate drifts outside a tolerance
band around the target, clamped to a configurable maximum drift from its
initial value. A verification run driving sixty consecutive
high-escalation-rate decisions through a policy instance with an attached
controller confirmed the threshold moves in the expected direction and
remains within its configured bound.

Both mechanisms default to their prior, static behaviour; the reputation
score defaults to zero and the adaptive controller is not attached unless
explicitly constructed, so no existing evaluation result from
Chapters~4--5 is affected by their introduction.

\section{Revisiting the Learning-Based Alternative}
\label{sec:app_c_dqn}

Given that this appendix concerns dynamic behavioural adjustment, it is
natural to ask whether the DQN baseline evaluated in \S5.2 --- or a
learning-based approach more generally --- would be a more practical
mechanism than the two described in \S4 of this appendix. The evidence
developed elsewhere in this thesis argues against it, for reasons
specific to this system rather than as a general claim about
reinforcement learning.

First, the DQN's own recorded training dynamics (\S5.2, Finding~5) show
it converging to a degenerate ``almost never ALLOW'' policy --- detection
rate rising from 87.6\% to 97.4\% within a single episode while
false-positive rate remained pinned near 1.0 --- as a direct consequence
of the dense-attack, single-intent training distribution used, rather
than as evidence of genuine threat discrimination. Any comparably-trained
deep reinforcement learning agent would be expected to reproduce this
failure mode unless the underlying training-distribution design were
corrected first, which would itself constitute substantial additional
work outside the scope of this appendix.

Second, this project's evaluation harness has an independently documented
history of subtle defects in exactly the kind of step-alignment logic a
reward computation would depend on: the project's bug log records two
separate instances in which a one-step misalignment between the
action-selecting observation and the label used to score it inflated
measured false-positive rates by up to an order of magnitude before
detection, in code that only needed to compute a descriptive statistic
after an episode had already concluded. A reward signal for a learning
agent is computed by structurally similar logic on every single
simulation step, and an equivalent undetected defect there would not
merely misreport a metric but would actively shape what the agent learns,
in a way that --- as with the original DQN's degenerate policy --- could
present as a coherent, high-performing solution while being fundamentally
unsound.

This reasoning is consistent with, and extends, the recommendation this
thesis already reaches in \S5.3 following \citet{rudin2019stop}: prefer
an interpretable model until a black-box alternative demonstrates clear,
validated superiority and can itself be explained after the fact. The
reputation override and the bounded threshold controller described in
\S4 of this appendix satisfy that standard directly --- both reduce to a
small number of auditable arithmetic operations rather than a trained
function --- and are accordingly proposed as the more practical near-term
mechanism for dynamic behavioural adjustment in this system.

Should a learning component nonetheless be desired in future work, the
most defensible form identified is a small contextual bandit --- not a
Q-network --- selecting among a handful of preset
\texttt{RATE\_THRESHOLD} values, trained on a reward derived from an
analyst's own confirmed true/false-positive label on a logged decision.
No such labelling pathway currently exists in this system; decisions are
logged in full (\texttt{DummyHoneypot}'s action audit log), but nothing
closes the loop from a human-confirmed judgement back into a trainable
signal. Constructing that pathway is identified as a prerequisite for any
future learning-based extension, rather than attempting to substitute a
proxy signal such as override trigger frequency, which would not
constitute genuine ground truth.

\section{Summary}

Table~\ref{tab:extensions_summary} summarises the six extensions
described in this appendix. All default to a configuration that
reproduces the prior exact behaviour of the system described in
Chapters~3--5, so the results reported in Chapter~4 remain valid
descriptions of the system's default configuration.

\begin{table}[ht]
  \centering
  \caption{Summary of extensions described in this appendix.}
  \label{tab:extensions_summary}
  \begin{tabular}{p{3.6cm}p{5.2cm}p{3.8cm}}
    \toprule
    Extension & Mechanism & Default state \\
    \midrule
    Session-coherent synthetic traffic & Per-episode intensity/persona profile & Always on \\
    Benign traffic diversity & Three weighted personas & Always on \\
    Randomised event payloads & Wordlist/jitter-range template resolution & Always on \\
    Escalation tracking & Severity-weighted EMA alongside the sliding window & Opt-in; window is default \\
    Cross-session reputation & Half-life-decayed offense score, new R4 override & Always tracked; R4 requires reputation $\geq 0.60$ \\
    Bounded threshold adaptation & Deadband controller on RATE\_THRESHOLD & Opt-in; static threshold is default \\
    \bottomrule
  \end{tabular}
\end{table}
